\chapter{Concepts}
\label{chap:foundations}

This chapter sets out the building blocks the prototype uses.
Each section defines the concept, states the property the
prototype relies on, and leaves the implementation detail for
Chapter~\ref{chap:system}. Section~\ref{sec:found-classical}
recaps the classical recommender families.
Section~\ref{sec:found-similarity} defines the similarity
measures and the time-decay function used throughout the system.
Section~\ref{sec:found-two-stage} states the two-stage
retrieve-and-rerank pattern. Section~\ref{sec:found-agents}
describes the LLM-agent loop with tool use and structured output.
Section~\ref{sec:found-grounding} develops hydration over
generation and the post-generation validation pipeline as a
structural defense against hallucination.
Section~\ref{sec:found-eval} defines the evaluation metrics used
to compare model and prompt configurations.

\section{Classical Recommender Foundations}
\label{sec:found-classical}

The classical layer uses the three families surveyed in
Chapter~\ref{chap:background}: item-based collaborative filtering
specialized as co-occurrence, content-based filtering over
weighted product attributes, and exponentially time-decayed
popularity for trending. The personalized slot uses a switching
hybrid as defined by Burke~\cite{Burke2002Hybrid}, picking the
algorithm per request based on how deep the user's purchase
history goes. Three properties of this layer matter for what
follows: every algorithm runs on the live catalog and returns
real product identifiers (it cannot invent products), every
algorithm is deterministic given the catalog state (the candidate
set is reproducible across evaluation runs), and every algorithm
runs inside a single Convex query within a tight latency budget
(Section~\ref{sec:found-two-stage}). Together they justify
treating the classical layer as the trusted source of candidates
and the LLM as a writer over those candidates.

\section{Similarity Measures and Time Decay}
\label{sec:found-similarity}

Four formal definitions support the classical layer.

\paragraph{Jaccard similarity.} Jaccard measures how much two
sets overlap, normalized by their combined size --- the share of
elements they share. For two finite sets $A$ and $B$,
\[
  J(A, B) \;=\; \frac{|A \cap B|}{|A \cup B|},
\]
with $J(\emptyset, \emptyset) := 0$ by convention. The coefficient
is symmetric, lives in $[0, 1]$, and equals one iff $A = B$.
Applied to the buyer sets of two products $p_a, p_b$, it produces
a similarity score that is robust to popularity differences: a
product with many buyers does not automatically score high
against another just because their intersection is
large~\cite{ComprehensiveRecSys2024}. The
prototype uses Jaccard over buyer sets for the ``frequently
bought together'' slot, with a minimum-support filter requiring
$|B(p_a) \cap B(p_b)| \ge 2$ for a pair to enter the precomputed
co-occurrence index. The filter suppresses spurious links from
incidental co-purchases.

\paragraph{Weighted attribute similarity.} Two products are
compared by averaging seven small attribute comparisons, with the
heaviest weight on the simplest signal (same category) and small
equal contributions from price, rating, brand, and the
LLM-enriched fields. Formally,
\[
  \mathrm{sim}(p_a, p_b) \;=\; \sum_{i=1}^{7} w_i \cdot f_i(p_a, p_b),
\]
where each $f_i$ is one of three primitives --- an indicator
$\mathbb{1}[\cdot]$ (1 if equal, 0 otherwise), the Jaccard
coefficient $J(\cdot, \cdot)$ over a tag-set argument, or a
proximity score $\mathrm{prox}(x, y) = 1 - |x-y|/\max(x,y)$ for
numeric attributes. Table~\ref{tab:found-sim-weights} lists the
components and their weights.

\begin{table}[ht]
  \centering
  \footnotesize
  \begin{tabular}{lll}
    \toprule
    Component & Weight & Primitive \\
    \midrule
    same category   & 0.25 & indicator \\
    tag overlap     & 0.20 & Jaccard \\
    use-case overlap & 0.15 & Jaccard \\
    good-for overlap & 0.10 & Jaccard \\
    same brand      & 0.10 & indicator \\
    price proximity & 0.10 & proximity \\
    rating proximity & 0.10 & proximity \\
    \bottomrule
  \end{tabular}
  \caption{Components of the weighted attribute similarity score.
    Weights sum to 1.0; category and tag overlap dominate.}
  \label{tab:found-sim-weights}
\end{table}

The \texttt{useCases} and \texttt{goodFor} fields carry implicit
attributes (a children's bicycle is also ``a gift for a
7-year-old'') and are produced once per product during seed by a
batch LLM call modeled on the COSMO knowledge graph~\cite{COSMO}.
At query time, $\mathrm{sim}$ is pure attribute arithmetic over
the denormalized product row.

\paragraph{Cosine similarity.} Cosine treats two vectors as
directions and measures the angle between them --- values near
one mean the vectors point the same way, values near zero mean
they are orthogonal. For two vectors $u, v \in \mathbb{R}^d$,
\[
  \cos(u, v) \;=\; \frac{\langle u, v \rangle}{\|u\| \cdot \|v\|}.
\]
The prototype uses cosine over text embeddings only inside the
review-theme-fidelity scorer of Section~\ref{sec:found-eval}:
each candidate review-summary theme is embedded and compared
against the embeddings of the source reviews, and themes whose
nearest review falls below a threshold $\tau$ are rejected.
Cosine beats Euclidean distance here because embedding norms
carry no semantic content.

\paragraph{Exponential time decay.} Trending is a weighted count
of recent events where each event's contribution shrinks
exponentially as it ages, so a sale from yesterday counts more
than a sale from a month ago. For a product $p$ with event set
$E_p$, per-event weight $w(\cdot)$ (a purchase weighs more than a
view), and event age $\mathrm{age}(e)$ measured in days from now,
\[
  \mathrm{trending}(p) \;=\; \sum_{e \in E_p} w(e) \cdot
  \exp\!\big(-\lambda \cdot \mathrm{age}(e)\big).
\]
The decay rate $\lambda$ is set by the desired half-life:
$\lambda = \ln 2 / T_{1/2}$. The prototype uses $T_{1/2} = 14$
days, giving $\lambda \approx 0.0495$. The score is precomputed
nightly into the \texttt{trendingScore} column on the
\texttt{products} table, so the trending slot is one indexed read
at query time. The half-life is the only meaningful knob, and it
depends on category; durable goods would justify a longer
half-life and fashion a shorter one~\cite{SessionRecSys2021}.

\section{The Two-Stage Retrieve-and-Rerank Pattern}
\label{sec:found-two-stage}

The two-stage retrieve-and-rerank pattern connects the classical
layer to the LLM layer. Its information-retrieval lineage runs
from cascade ranking in web search~\cite{HSTU2024} through to
production e-commerce deployments~\cite{ShopifyHSTU}. The thesis
adopts it for a specific reason: it is the structural answer to
the question Chapter~\ref{chap:background} raised about Rufus,
namely how an LLM-assisted shopping system avoids hallucinating
products and specifications.

Formally, let $\mathcal{C}$ be the full product catalog and $q$ a
slot intent or user query. A retrieve function $R$ produces a
small candidate set
\[
  S \;=\; R(q, \mathcal{C}), \qquad |S| \ll |\mathcal{C}|,
  \qquad S \subseteq \mathcal{C},
\]
under a tight latency budget $\tau_R$. A rerank-and-explain
function $L$ then runs over $S$ alone:
\[
  (\hat{S}, r) \;=\; L(q, S, \textrm{ctx}), \qquad \hat{S} \subseteq S,
\]
producing a final ranked subset $\hat{S}$ and a free-text
rationale $r$. The LLM never reads $\mathcal{C}$ directly. By
construction $\hat{S} \subseteq S \subseteq \mathcal{C}$, so every
product identifier reaching the user is a real catalog item.

Three properties of the pattern matter for trust and cost.
First, the inclusion $\hat{S} \subseteq \mathcal{C}$ is enforced
structurally, not by prompt: identifiers the LLM emits that are
not in $S$ are dropped at the validation layer
(Section~\ref{sec:found-grounding}). Second, the latency budget
is split: $\tau_R$ is the classical query budget (the prototype
targets $\tau_R < 50$ ms at p95), and the rest of the
consumer-acceptable end-to-end budget goes to $L$. Third, LLM
inference cost scales with input context size, so bounding $|S|$
caps the per-request token cost regardless of
$|\mathcal{C}|$~\cite{HSTU2024}. At the prototype scale
($|\mathcal{C}| \approx 200$) the second property is comfortable;
the bound on $|S|$ is what keeps the pattern viable as
$|\mathcal{C}|$ grows.

The pattern also splits responsibilities cleanly. The retrieve
function $R$ owns relevance and recall over the full catalog and
only ever sees structural data. The rerank-and-explain function
$L$ owns natural-language presentation, conflict surfacing, and
tradeoff articulation, and only ever sees the candidate set plus
context. Neither concern leaks into the other.

\section{LLM Agents and Tool Use}
\label{sec:found-agents}

In the prototype, $L$ is an LLM \emph{agent}: a loop in which the
model alternates between reasoning, picking a tool, executing it,
and reading the observations it returns. The pattern is the
ReAct loop, introduced for multi-step reasoning over external
data~\cite{ReAct2022}, and now standard in tool-augmented LLM
deployments. A single iteration produces a triple
\[
  (\mathrm{thought}_t, \mathrm{action}_t, \mathrm{observation}_t),
\]
where $\mathrm{action}_t$ is either a tool call or a terminal
text response, and $\mathrm{observation}_t$ is the tool's
structured return value. The loop ends when the model emits a
final text response or when the iteration count hits a configured
step budget $B$.

Three engineering choices follow from this loop.

\paragraph{Tool calling.} The agent gets a typed tool catalog at
the start of every turn. Each tool has a name, a JSON schema for
its arguments, and a deterministic implementation that the agent
runtime executes in a sandbox with access to the catalog. The
prototype exposes seven tools to the advisor:
\texttt{getProductDetails}, \texttt{searchProducts},
\texttt{getRecommendations}, \texttt{getCartContents},
\texttt{getReviewSummary}, \texttt{compareProducts}, and
\texttt{getCategoryStructure}. Each tool is a thin wrapper over a
Convex query, so the agent reads the same data the production
client reads. The tool catalog is the only way the LLM can pull
information from the live catalog during an advisor turn.

\paragraph{Structured output.} For surfaces that render
structured UI (the productId list on the advisor surface, the
comparison table, the review-summary object), the final agent
response has to validate against a JSON schema. Native JSON-mode
support in modern LLMs~\cite{StructuredOutput2024} pushes the
adherence rate above 99\% for moderate-depth schemas, so the
prototype treats schema validation as a hard contract rather than
a best-effort guardrail. A schema-invalid response counts as a
failed turn and is not surfaced to the user.

\paragraph{Step budget and termination.} The step budget $B$ caps
the loop length to bound worst-case latency and cost. The
prototype uses $B = 12$, set after an early bug class where the
agent ran out of steps making redundant tool calls and never
produced a terminal text response. The eval harness
(Section~\ref{sec:found-eval}) tracks both the
\texttt{hasFinalAnswer} property (did the agent terminate before
$B$?) and the \texttt{toolCallEfficiency} property (how many of
the calls were duplicates?), and any regression on either shows
up immediately in the configuration sweep.

The agent loop's cost is dominated by the LLM call at each
iteration. A turn with three tool calls and a terminal response
means four LLM calls. The classical layer's tight latency budget
(Section~\ref{sec:found-two-stage}) is what makes this loop
viable: each tool call resolves in milliseconds, so the total
wall-clock time is dominated by LLM inference, not tool
execution.

\section{Grounded Generation and Output Validation}
\label{sec:found-grounding}

Hallucination is the LLM behavior where the model produces
plausible-sounding output that is not supported by its inputs. In
a shopping context, the failure modes catalogued in
Chapter~\ref{chap:background} are all instances: invented product
identifiers, fabricated prices, mismatched specifications, and
review themes that do not appear in the source corpus. The
generative behavior that produces these failures is baked into
the model class~\cite{HallucinationSurvey2023}; mitigation has to
be structural.

\paragraph{Hydration over generation.} The prototype enforces a
single principle: the LLM emits structured references, and the
client renders live data from the catalog. Formally, let $\hat{S}$
be the set of productIds the agent emits and let $H : \hat{S} \to
\mathcal{P}$ be the hydration function that maps each identifier
to the canonical product record fetched at render time from the
live catalog. The user-visible surface is
\[
  \{ H(p) \,:\, p \in \hat{S} \} \;\cup\; \mathrm{sanitize}(r),
\]
where $r$ is the agent's free-text rationale and $\mathrm{sanitize}$
removes price and specification claims (defined below). The LLM
never sees a price, never produces a price for display, and never
renders a product attribute. Every factual surface the user sees
was read from the catalog at render time.

\paragraph{Validation pipeline.} The hydration principle is
backed by a post-generation validation pipeline with three
sequential checks. Let $\mathcal{C}_{id}$ denote the set of
product identifiers in the live catalog at the moment of the
agent turn.

\begin{enumerate}
  \item \textbf{Identifier-existence.} The structured part of the
    response is intersected with $\mathcal{C}_{id}$:
    \[
      \hat{S}' \;=\; \hat{S} \cap \mathcal{C}_{id}.
    \]
    Identifiers in $\hat{S} \setminus \mathcal{C}_{id}$ are dropped
    silently. The check is a single indexed read per identifier and
    runs inside the same Convex action that wraps the agent call,
    so the client never sees an unverified identifier. This is the
    primary defense against the 32\% recommendation accuracy
    failure mode documented for
    Rufus~\cite{RufusWrong2024}.
  \item \textbf{Claim stripping.} The free-text rationale $r$ is
    run through a deterministic regex pass that strips tokens
    matching price patterns (currency-prefixed numbers, decimals
    next to keywords like ``costs'', ``priced at'', ``retails
    for'') and specification patterns (numbers next to unit
    keywords like ``GB'', ``inches'', ``mAh'', ``Hz'', ``W'').
    The sanitized output is $\mathrm{sanitize}(r)$. The system
    prompt forbids the model from quoting these in the first
    place; stripping is a belt-and-braces measure that catches
    prompt-following lapses~\cite{ShopifyAgentic}.
  \item \textbf{Theme fidelity (batch path).} For review summaries
    only, each claimed theme $t$ is checked against the embedding
    set $\mathcal{E}_{\mathrm{rev}}$ of the source review corpus by
    a nearest-neighbor lookup under cosine similarity
    (Section~\ref{sec:found-similarity}). A theme is accepted iff
    \[
      \max_{e \in \mathcal{E}_{\mathrm{rev}}} \cos(\mathrm{embed}(t), e)
      \;\geq\; \tau,
    \]
    with threshold $\tau$ tuned on a small held-out batch. Themes
    that fail the test are dropped before the summary is
    persisted. This is the defense against the review-fabrication
    failure mode~\cite{RufusConfidentlyWrong2024}.
\end{enumerate}

The three checks are not redundant. Identifier-existence catches
structural hallucinations. Claim stripping catches drift in free
text. Theme fidelity catches fabrication in structured outputs
that pass the schema check but are not supported by source data.
Together they encode the design constraint that the LLM is a
candidate selector and a writer, never a source of facts.

\section{Evaluation Metrics for Hybrid Systems}
\label{sec:found-eval}

A hybrid system with a classical layer and an LLM layer cannot be
evaluated by a single metric. The prototype's evaluation rests on
three layers of scorer, each addressing a different failure mode.

\paragraph{Programmatic scorers.} A programmatic scorer is a
deterministic function from the agent's response and the live
catalog snapshot to a boolean or a real in $[0, 1]$. Five scorers
are used in the canonical sweep.

\begin{itemize}
  \item \texttt{hasFinalAnswer}: did the agent terminate before the
    step budget? Catches the maxSteps-cutoff bug class.
  \item \texttt{groundedness}: is every productId cited in the
    answer in the catalog?
    Formally, $|\hat{S} \setminus \mathcal{C}_{id}| / |\hat{S}|$,
    inverted so higher is better.
  \item \texttt{factuality}: do claimed prices and ratings match
    the database snapshot captured at run start? Scored as the
    fraction of claims that match.
  \item \texttt{reviewThemeFidelity}: do claimed review themes
    have an embedding match in the source corpus above
    threshold $\tau$? Generalizes the validator described in
    Section~\ref{sec:found-grounding}.
  \item \texttt{toolCallEfficiency}: a duplicate-call penalty
    plus the fraction of the step budget used.
\end{itemize}

Programmatic scorers are cheap, reproducible, and free of model
bias. They are the harness's strongest signal.

\paragraph{LLM-as-judge scorers.} Properties that cannot be scored
deterministically (completeness, helpfulness, tradeoff surfacing,
tool appropriateness, tone) are scored by a separate judge LLM on
a five-point rubric. Two design choices reduce judge bias. First,
the judge is from a different model family than any configuration
under test: the prototype uses Claude Haiku 4.5 as the judge for
sweeps containing GPT and Gemini configurations, which avoids the
well-documented self-preference bias of within-family
judging~\cite{LLMJudgeSurvey2024}. Second, the judge is given the
question and the answer only, not the configuration identifier,
which prevents prestige bias.

\paragraph{Pareto and composite ranking.} A configuration $c$ is
described by a vector of measurements: average judge quality
$Q(c)$, average programmatic correctness $C(c)$, per-query cost
$\mathrm{cost}(c)$, p95 latency $\mathrm{lat}(c)$, and tool-step
efficiency $E(c)$. The sweep produces this vector for every
configuration. Two reporting forms are used.

The pareto front on $(\mathrm{cost}, Q)$ identifies the
configurations that no other configuration strictly dominates on
both axes: for any $c$ on the front, no $c'$ exists with
$\mathrm{cost}(c') < \mathrm{cost}(c)$ and $Q(c') > Q(c)$. The
same construction yields the pareto front on
$(\mathrm{lat}, Q)$. The front is the natural object for choosing
a deployment configuration under explicit budget.

A composite score collapses the measurement vector into a single
number by taking a weighted average across the five axes, used
when one ranked list is needed. The prototype defines
\[
  \mathrm{composite}(c) \;=\;
  0.30 \cdot Q(c) +
  0.25 \cdot C(c) +
  0.20 \cdot (1 - P_{\mathrm{cost}}(c)) +
  0.15 \cdot (1 - P_{\mathrm{lat}}(c)) +
  0.10 \cdot E(c),
\]
where $P_{\mathrm{cost}}(c)$ and $P_{\mathrm{lat}}(c)$ are the
cost and latency of $c$ normalized linearly into $[0, 1]$ against
the cheapest and fastest configurations in the sweep. The weights
are an explicit editorial choice: quality dominates, correctness
is the second largest factor, cost and latency contribute
moderately, and efficiency is a tie-breaker. The composite is
always reported alongside the pareto front, never instead, so a
reader who disagrees with the weights can read the underlying
tradeoffs directly.

The metrics defined here are the same ones reported in the
results section of Chapter~\ref{chap:system}. The next chapter
moves from the abstract architecture to the concrete prototype:
the codebase, the data layer, the classical algorithm
implementations, the agent infrastructure, the validation
pipeline, the frontend surfaces, the evaluation harness, and the
measured numbers across the canonical configuration sweep.
